% Removed 2026-10-19 (cut round two): full previous text of sections/appC_square.tex before cutting.
\section{Square grids}\label{app:square}

This section answers question~(4) of Section~\ref{sec:open} from below. Let $n\ge2$, let $\mathcal R$ be the set of rectangles of the $n\times n$ grid, tilted ones included, and let $L=1+\ln n$. For a point $v$, let $d_v$ be the number of rectangles with corner $v$. By Lemma~\ref{lem:threecorners}, two different rectangles share at most two corners.

\begin{lemma}\label{lem:compress}
Let $R_1\ne R_2$ be rectangles with exactly two common corners $p$ and $q$. If $\{p,q\}$ is a side of both, or a diagonal of both, then $C(R_1)\triangle C(R_2)$ is the corner set of a rectangle.
\end{lemma}

\begin{proof}
If $\{p,q\}$ is a side of both, the other corners are $p+u,q+u$ and $p+v,q+v$ with $u\ne v$ nonzero and perpendicular to $q-p$. They are distinct, as $p+u=q+v$ would make $v-u=p-q$ perpendicular to itself, and form a parallelogram with perpendicular sides $q-p$ and $v-u$. If $\{p,q\}$ is a diagonal of both, the other corners are $o\pm u$ and $o\pm v$ with $o=(p+q)/2$, $|u|=|v|$ and $v\ne\pm u$; two segments with a common midpoint and equal length are the diagonals of a rectangle.
\end{proof}

Let $r_2(t)$ be the number of integer pairs $(a,b)$ with $a^2+b^2=t$, and $\tau(t)$ the number of divisors of $t$. Hardy and Wright \cite[Sections~16.9, 16.10, and~18.1]{hardywright2008} prove that $r_2(t)=4(d_1(t)-d_3(t))\le4\tau(t)$, where $d_i(t)$ counts the divisors of $t$ congruent to $i$ modulo $4$, and that for each $\varepsilon>0$ there is $K_\varepsilon$ with $\tau(t)\le K_\varepsilon t^\varepsilon$ for all $t\ge1$.

Let $\mathcal A$ be the family of all corner sets $C(R)$ with $R\in\mathcal R$, together with all sets $C(R_1)\triangle C(R_2)$ for pairs $R_1\ne R_2$ whose exactly two common corners form a side of one and a diagonal of the other. Let $\mathcal B$ be the family of all sets $C(R_1)\triangle C(R_2)$ for pairs $R_1\ne R_2$ with exactly one common corner.

\begin{lemma}\label{lem:counts}
\begin{enumerate}[(a)]
\item For every point $v$, $d_v\le8n^2L$, and $|\mathcal R|\le2n^4L$.
\item For every $\varepsilon>0$, $|\mathcal A|\le A_\varepsilon n^{4+2\varepsilon}L$ with $A_\varepsilon=2(1+16\cdot2^\varepsilon K_\varepsilon)$, and $|\mathcal B|\le32n^6L^2$.
\item A blackout that contains no member of $\mathcal A\cup\mathcal B$ is valid.
\end{enumerate}
\end{lemma}

\begin{proof}
(a) The sides of a rectangle at $v$ are $k(a,b)$ and $l(-b,a)$ with $\gcd(a,b)=1$, $k\ge1$, and $l\ne0$ integers. With $d=\max(|a|,|b|)$ there are $8d$ choices of $(a,b)$, at most $(n-1)/d$ of $k$, and at most $2(n-1)/d$ of $l$. Each rectangle gives two ordered pairs of sides, so $2d_v\le\sum_{d=1}^{n-1}16(n-1)^2/d\le16n^2L$. Then $4|\mathcal R|=\sum_vd_v\le8n^4L$.

(b) Fix $R_1$ and a side $\{p,q\}$ of it. A rectangle $R_2$ with diagonal $\{p,q\}$ is determined by its other diagonal $p'-q'$, an integer vector with $|p'-q'|^2=|p-q|^2<2n^2$, so there are at most $4K_\varepsilon(2n^2)^\varepsilon$ such $R_2$. Hence $|\mathcal A|\le|\mathcal R|(1+16K_\varepsilon2^\varepsilon n^{2\varepsilon})\le A_\varepsilon n^{4+2\varepsilon}L$. Each member of $\mathcal B$ comes from two rectangles at a common corner $v$, so $|\mathcal B|\le\sum_v\binom{d_v}2\le\frac{n^2}2(8n^2L)^2$.

(c) If a blackout $S$ is not valid, some $R_1\ne R_2$ have the same presentation, and $C(R_1)\triangle C(R_2)\subseteq S$ by Lemma~\ref{lem:diff}. With two common corners this set lies in $\mathcal A$, by Lemma~\ref{lem:compress} or by the definition of $\mathcal A$; with one it lies in $\mathcal B$, and with none $C(R_1)\subseteq S$ lies in $\mathcal A$.
\end{proof}

\begin{theorem}[random blackouts with deletions]\label{thm:alteration}
We have
\[
\liminf_{n\to\infty}\frac{f(n,n)\,(1+\ln n)^{2/5}}{n^{6/5}}\ge\frac5{6\cdot192^{1/5}}=0.2911\ldots,
\]
so for every $c<5/(6\cdot192^{1/5})$ we have $f(n,n)\ge c\,n^{6/5}(1+\ln n)^{-2/5}$ for all large $n$. In particular $f(n,n)-2n\to\infty$.
\end{theorem}

\begin{proof}
Following Alon and Spencer~\cite[Section~3.2]{alonspencer2016}, put $a=192^{-1/5}$, $T=n^{6/5}L^{-2/5}$, and $\theta=a\,n^{-4/5}L^{-2/5}\in(0,1)$, and select each grid point independently with probability $\theta$. Let $X$, $Y$, and $Z$ count the selected points and the fully selected members of $\mathcal A$ and of $\mathcal B$. Then
\[
\mathbb E(X-Y-Z)=n^2\theta-|\mathcal A|\theta^4-|\mathcal B|\theta^6 .
\]
Here $n^2\theta=aT$, and Lemma~\ref{lem:counts}(b) with $\varepsilon=1/10$ gives $|\mathcal B|\theta^6\le32a^6T$ and $|\mathcal A|\theta^4\le A_{1/10}a^4n^{-1/5}T$. Since $32a^5=\frac16$, some selection $S_0$ has $X-Y-Z\ge(\frac56a-A_{1/10}a^4n^{-1/5})T$. Deleting one point from each member of $\mathcal A\cup\mathcal B$ inside $S_0$ leaves a set $S$ with $|S|\ge X-Y-Z$, valid by Lemma~\ref{lem:counts}(c). Hence $f(n,n)/T\ge\frac56a-A_{1/10}a^4n^{-1/5}$, which tends to $5/(6\cdot192^{1/5})$. Finally $T/n\to\infty$, which gives $f(n,n)-2n\to\infty$.
\end{proof}

% Round 3 (2026-10-19): sentence cut from the statement of Theorem thm:alteration (restates the liminf).
so for every $c<5/(6\cdot192^{1/5})$ we have $f(n,n)\ge c\,n^{6/5}(1+\ln n)^{-2/5}$ for all large $n$.
